\section{Discussion}\label{sec:07}
On the spherical history the two round boundary metrics determine the same instantaneous spectral subspaces. The failure of their implementability under propagation depends on the first time derivative of shape. Including this derivative at the endpoints removes the leading mismatch. The next surviving boundary term has nonzero coefficient on a band of magnetic labels; oscillatory phases suppress some individual modes but leave a positive band average. This explains why the same first-order prescription yields finite number and divergent energy. The compact Nil example exhibits a different lower-bound mechanism: one endpoint is shear-free, so no two-end interference has to be averaged.

Both examples start with a prescribed external field, geometry and family of particle definitions. Their conclusions give concrete criteria for finite endpoint observables on these backgrounds. Extending the energy threshold to massive fields or to general Nil histories requires new high-mode estimates. A local stress tensor would additionally involve state regularity and covariant subtraction. On a generic anisotropic three-manifold the number of independent labels and the spectral multiplicities may differ; these enter explicitly in both proofs.
